\documentclass[10pt,a4paper]{amsart}
\usepackage[margin=19mm]{geometry}
\usepackage{amsmath,amssymb,mathtools}
\usepackage[scr=rsfs,cal=euler]{mathalfa}
\usepackage{xfrac}
\usepackage[unicode,hidelinks,bookmarks=false]{hyperref}
\newcommand{\ie}{i.e.\ }
\newcommand{\cf}{cf.\ }
\newcommand{\resp}{respectively\ }
\newcommand{\Subsection}{\subsection}
\providecommand{\nobreakdash}{\nobreak\hbox{-}}
\setlength{\emergencystretch}{2em}
\multlinegap0pt
\usepackage{bm}
\usepackage{bbm}
\usepackage{dsfont}
\usepackage{xspace}
\usepackage{cancel}
\usepackage{mathtools}
\usepackage{amsthm}
\makeatletter
\@ifundefined{theorem}{\newtheorem{theorem}{Theorem}}{}
\@ifundefined{lemma}{\newtheorem{lemma}[theorem]{Lemma}}{}
\@ifundefined{proposition}{\newtheorem{proposition}[theorem]{Proposition}}{}
\makeatother
\makeatletter
\newcommand{\assign}{\coloneqq}
\expandafter\ifx\csname descriptioncompact\endcsname\relax
\newenvironment{descriptioncompact}{\begin{description} }{\end{description}}
\fi
\expandafter\ifx\csname enumeratealpha\endcsname\relax
\newenvironment{enumeratealpha}{\begin{enumerate}[a{\textup{)}}] }{\end{enumerate}}
\fi
\expandafter\ifx\csname enumerateroman\endcsname\relax
\newenvironment{enumerateroman}{\begin{enumerate}[i.] }{\end{enumerate}}
\fi
\expandafter\ifx\csname itemizedot\endcsname\relax
\newenvironment{itemizedot}{\begin{itemize} \renewcommand{\labelitemi}{$\bullet$}\renewcommand{\labelitemii}{$\bullet$}\renewcommand{\labelitemiii}{$\bullet$}\renewcommand{\labelitemiv}{$\bullet$}}{\end{itemize}}
\fi
\expandafter\ifx\csname proof\endcsname\relax
\newenvironment{proof}{\noindent\textbf{Proof\ }}{\hspace*{\fill}$\Box$\medskip}
\fi
\newcommand{\tmop}[1]{\ensuremath{\operatorname{#1}}}
\makeatother
\title[Taylor expansions over generalised power series]{Taylor expansions over generalised power series}
\author{Vincent Bagayoko, Vincenzo Mantova}
\date{Source: arXiv:2509.08473v1 (CC BY 4.0)}
\begin{document}
\maketitle

\noindent\emph{Source and licence.} Taylor expansions over generalised power series. Version arXiv:2509.08473v1 (CC BY 4.0), distributed under the Creative Commons Attribution 4.0 International licence, as recorded in the arXiv licence metadata for this exact version and shown on the abstract page https://arxiv.org/abs/2509.08473v1. Full rights notice: licenses/arxiv-2509.08473-CC-BY-4.0.txt.\par\smallskip

\noindent\textit{Editorial scope.} Counted unit: the Proposition whose source label is \texttt{prop-Noeth-ext-der} in the article (source0.tex lines 2105--2107, source label \texttt{prop-Noeth-ext-der}), delivered together with the article's own complete original proof of it (source0.tex lines 2109--2183, one \texttt{proof} environment of 75 lines). No mathematics is written, repaired or re-derived: statement and proof are the article's own text line for line. Required context retained uncounted: lem:coinitial-sequence (lines 2036--2046). Every definition, display and label of those lines is retained unchanged. Omitted from this package and not redistributed: 1--2035, 2047--2104, 2108--2108, 2184--2888. Nothing inside a retained excerpt is omitted or altered: every retained body is delivered exactly as the article's lines are written, so no omission receipt of any kind accompanies this package. Citations: the retained excerpts carry no citation command at all, so no bibliography entry of the article is transcribed and no reference is redistributed. Reference adaptation: none was needed and none was made; every reference inside the delivered unit resolves inside it, so no unresolved reference or citation is produced. Printed slips kept literal: no printed word, symbol, hypothesis or number of the counted statement or of its proof was altered. Layout and class: the article's own class is \texttt{article} with packages including fontenc, babel, hyperref, authblk, amsmath, amssymb, amsthm, graphicx, stmaryrd, enumerate, cleveref, latexsym, makeidx, mathtools; the wrapper is the lane's pinned \texttt{amsart} editorial wrapper at 10pt on A4, which restates the theorem-like environments the retained text needs and adds the article's own macro definitions that the retained text needs (\texttt{\textbackslash assign}, \texttt{\textbackslash tmop}), with extra wrapper packages (bm, bbm, dsfont, xspace, cancel, mathtools). The delivered numbering is the unit's own. Assets: no retained span contains a figure environment, a graphics-inclusion command, a PDF-page inclusion, a table or a TikZ picture; no image file is copied into this package, the package declares no asset dependency and no image file is redistributed with it. Licence: version arXiv:2509.08473v1 of this article is distributed under the Creative Commons Attribution 4.0 International licence, as recorded in the arXiv licence metadata for this exact version and shown on the abstract page https://arxiv.org/abs/2509.08473v1. A full rights notice is delivered at licenses/arxiv-2509.08473-CC-BY-4.0.txt. Characters: every retained excerpt is pure ASCII, so the delivered engine, XeLaTeX with its default Latin Modern text font, reports no missing character.
\begin{lemma}
  \label{lem:coinitial-sequence}Let $C$ be a convex subset of some truncation
  closed $L \subseteq \mathbb{S}$. Then there is a series $\varphi$ and a
  coinitial sequence in $C$ of the form $(\psi_i + c_i \mathfrak{o}_i)_{i <
  \kappa}$, where each $\psi_i$ is a truncation of $\varphi$, $c_i \in
  \mathbf{K}$, $\mathfrak{o}_i \in \mathfrak{M} \cup \{ 0 \}$. Moreover, we
  may assume that either $(\psi_i)_{i < \kappa}$ is injective, in which case
  $\mathfrak{o}_i \in \tmop{supp} (\varphi) \cup \{ 0 \}$ (in particular,
  $(\psi_i + r_i \mathfrak{o}_i)_{i < \kappa}$ is weakly summable), or $\psi_i
  = \varphi$ for all $i < \kappa$.
\end{lemma}

\begin{proposition}
  \label{prop-Noeth-ext-der}The function $\overline{\partial}$ is Noetherian.
\end{proposition}

\begin{proof}
  Let $(\mathfrak{m}_i X^{k_i})_{i \in \mathbb{N}}$ be a strictly
  $\prec_{\mathfrak{S}}$-decreasing sequence in $\mathfrak{M} \times
  X^{\mathbb{N}}$. This means, by definition, that $(k_i)_{i \in \mathbb{N}}$
  is weakly increasing and that there are monomials $\mathfrak{u}_i \succ
  \mathfrak{S}$ such that $\mathfrak{m}_{i + 1} \mathfrak{u}_i^{k_{i + 1} -
  k_i} \prec \mathfrak{m}_i$. Letting $\mathfrak{p}_{i + 1} \assign
  \mathfrak{u}_0^{k_1 - k_0} \cdots \mathfrak{u}_i^{k_{i + 1} - k_i}$,
  $\mathfrak{p}_0 \assign 1$, we find that $\mathfrak{m}_{i + 1}
  \mathfrak{p}_{i + 1} \prec \mathfrak{m}_i \mathfrak{p}_i$.
  
  Now pick some arbitrary $\mathfrak{n}_i \in \tmop{supp} \mathfrak{m}_i'$ for
  $i \in \mathbb{N}$. We claim that after taking a subsequence, the monomials
  $\mathfrak{n}_i \mathfrak{p}_i$ appear in the supports of some summable
  family. This implies that there are $i < j$ such that $\overset{}{}
  \mathfrak{n}_i \mathfrak{p}_i \succ \mathfrak{n}_j \mathfrak{p}_j$, and so
  $\mathfrak{n}_i \succ (\min_{i \leqslant n < j} \mathfrak{u}_n) ^{k_j - k_i}
  \mathfrak{n}_j$, thus $\mathfrak{n}_i X^{k_i} \succ_{\mathfrak{S}}
  \mathfrak{n}_j X^{k_j}$, proving that $(\mathfrak{m}_i' X^{k_i})_{i \in
  \mathbb{N}}$ is summable, and so that $\overline{\partial}$ is Noetherian.
  
  As a warm-up, observe that if $(k_i)_{i \in \mathbb{N}}$ is constant, then
  $\mathfrak{n}_i \mathfrak{p}_i =\mathfrak{n}_i$ is in the support of
  $\mathfrak{m}_i'$, and $(\mathfrak{m}_i')_{i \in \mathbb{N}}$ is summable by
  strong linearity of $\partial$.
  
  In the general case, observe that by strong linearity of $\partial$, the
  family
  \[ (\mathfrak{m}_i \mathfrak{p}_i)' =\mathfrak{m}_i' \mathfrak{p}_i
     +\mathfrak{m}_i \mathfrak{p}_i' =\mathfrak{p}_i  (\mathfrak{m}_i'
     +\mathfrak{m}_i \mathfrak{p}_i^{\dag}) \]
  is summable. Note moreover that $\mathfrak{n}_i \mathfrak{p}_i \in
  \tmop{supp} \mathfrak{m}_i' \mathfrak{p}_i$. We also have
  \[ \mathfrak{p}_{i + 1}^{\dagger} = (k_1 - k_0) \mathfrak{u}_0^{\dagger} +
     \cdots + (k_{i + 1} - k_i) \mathfrak{u}_i^{\dag} . \]
  If the sequence $(\ell (\mathfrak{u}_i))_{i \in \mathbb{N}}$ is weakly
  summable, then $(\mathfrak{u}_i^{\dagger})_{i \in \mathbb{N}}$ is weakly
  summable, so $ (\mathfrak{m}_i \mathfrak{p}_i')_{i \in \mathbb{N}} =
  ((\mathfrak{m}_i \mathfrak{p}_i) \mathfrak{p}_i^{\dagger})_{i \in
  \mathbb{N}}$ is summable, hence $(\mathfrak{m}_i' \mathfrak{p}_i)_{i \in
  \mathbb{N}}$ is summable, and we are done.
  
  If the sequence $(\ell (\mathfrak{u}_i))_{i \in \mathbb{N}}$ is not
  coinitial in $\ell (\mathfrak{M} \setminus \mathfrak{S})$, then there is
  $\mathfrak{u} \succ \mathfrak{S}$ such that $\mathfrak{u} \preccurlyeq
  \mathfrak{u}_i$ for all $i$. So we may assume that $\mathfrak{u}_i
  =\mathfrak{u}$ for all $i$, in which case $(\ell (\mathfrak{u}_i))_{i \in
  \mathbb{N}}$ is weakly summable, and we are done. Otherwise, we may replace
  $(\mathfrak{u}_i)_{i \in \mathbb{N}}$ with a subsequence of any other
  coinitial sequence in $\ell (\mathfrak{M} \setminus \mathfrak{S})$. Since
  $\ell (\mathfrak{M})$ is truncation closed and $\ell (\mathfrak{M} \setminus
  \mathfrak{S})$ is a convex subset of $\ell (\mathfrak{M})$, by
  Lemma~\ref{lem:coinitial-sequence}, we may choose the sequence so that $(\ell
  (\mathfrak{u}_i))_{i \in \mathbb{N}}$ is either weakly summable, or of the
  form $(\varphi + c_i \mathfrak{o}_i)_{i \in \mathbb{N}}$ and not weakly
  summable. Thus after taking a subsequence with $(\mathfrak{o}_i)_{i \in
  \mathbb{N}}$ strictly increasing and with $r_i < 0$.
  
  In the former case, we are done. In the latter, note that for any choice of
  non-zero $n_i \in \mathbb{N}$, there is another coinitial sequence
  $(\mathfrak{v}_i)_{i \in \mathbb{N}}$ in $\ell (\mathfrak{M} \setminus
  \mathfrak{S})$ such that $\ell (\mathfrak{v}_i) = \varphi + n_i c_i
  \mathfrak{o}_i$. If there are $j \in \mathbb{N}$ and infinitely many $i \in
  \mathbb{N}$ such that $\mathfrak{n}_i \in \tmop{supp} (\mathfrak{m}_i
  \varphi') \cup \tmop{supp} (\mathfrak{m}_i \mathfrak{o}_j')$, we note that
  $(\mathfrak{m}_i \mathfrak{p}_i  (\varphi' +\mathfrak{o}_j'))_{i \in
  \mathbb{N}}$ is summable, and we are done. Otherwise, we can choose $n_i$ so
  that $\mathfrak{n}_i \in \tmop{supp} (\mathfrak{m}_i \mathfrak{o}_j')$
  implies that $\mathfrak{n}_i \in \tmop{supp} (\mathfrak{m}_i'
  +\mathfrak{m}_i \mathfrak{p}_i^{\dag})$ and in particular $\mathfrak{n}_i
  \mathfrak{p}_i \in \tmop{supp} (\mathfrak{m}_i \mathfrak{p}_i)'$. With this
  choice, for each $i$ we have $\mathfrak{n}_i \mathfrak{p}_i \in (\tmop{supp}
  \mathfrak{m}_i \varphi') \cup (\tmop{supp} (\mathfrak{m}_i
  \mathfrak{p}_i)')$, and we are done.
\end{proof}

\end{document}
